% !TEX root = ../main.tex
\begin{ouverture}[flamme]{10}{Combustions et énergie chimique}{Partie 1 · Constitution et transformations de la matière}
{Combustible\enspace·\enspace Comburant\enspace·\enspace Énergie de réaction\enspace·\enspace Pouvoir calorifique\enspace·\enspace Énergie de liaison}
\objectif{Citer des combustibles usuels et modéliser une combustion par une réaction d'oxydoréduction.}
\objectif{Écrire l'équation de la combustion complète d'un alcane et d'un alcool.}
\objectif{Utiliser l'énergie molaire de réaction et le pouvoir calorifique pour calculer une énergie libérée.}
\objectif{Estimer une énergie de réaction en phase gazeuse à partir des énergies de liaison.}
\objectif{Citer des applications des combustions, leurs risques et des pistes de développement durable.}
\prerequis{Réactions d'oxydoréduction (chapitre~3) ; avancement (chapitre~4).}
\prerequis{Schémas de Lewis (chapitre~6) ; alcanes et alcools (chapitre~8).}
\end{ouverture}
\essentiel
\essSection{1}{Combustibles et combustions}
\begin{definition}[Combustion]
Une \textbf{combustion} est une transformation chimique exothermique entre un \textbf{combustible} et un \textbf{comburant}, le
plus souvent le dioxygène de l'air. C'est une réaction d'oxydoréduction : le combustible est le réducteur, le dioxygène
l'oxydant. Les combustibles usuels sont d'origine fossile (gaz naturel, butane, propane, essence, gazole, fioul, charbon)
ou renouvelable (bois, biogaz, bioéthanol).
\end{definition}
\schema{10-triangle}
\essSection{2}{Écrire l'équation d'une combustion}
\begin{propriete}[Combustion complète]
La combustion complète d'une espèce formée de carbone, d'hydrogène et éventuellement d'oxygène produit du dioxyde de
carbone et de l'eau. Pour un alcane et pour un alcool :
\[\mathrm{C}_n\mathrm{H}_{2n+2}+\frac{3n+1}{2}\,\ce{O2}\longrightarrow n\,\ce{CO2}+(n+1)\,\ce{H2O}\qquad\qquad\mathrm{C}_n\mathrm{H}_{2n+2}\mathrm{O}+\frac{3n}{2}\,\ce{O2}\longrightarrow n\,\ce{CO2}+(n+1)\,\ce{H2O}.\]
On ajuste C, puis H, puis O, en multipliant par 2 si nécessaire : $\ce{2C4H10 + 13O2 -> 8CO2 + 10H2O}$. Si le dioxygène manque,
la combustion est \textbf{incomplète} : elle produit aussi du monoxyde de carbone $\ce{CO}$, toxique, ou du carbone (suie).
\end{propriete}
\essSection{3}{L'énergie libérée}
\begin{definition}[Énergie molaire de réaction, pouvoir calorifique]
L'\textbf{énergie molaire de réaction} $E_\mathrm{r}$, en \unit{kJ.mol^{-1}}, est l'énergie échangée par le système pour un avancement de
\qty{1}{mol} ; elle est négative pour une réaction exothermique. La combustion de $n$ moles de combustible (nombre stœchiométrique 1)
libère $Q=n\times|E_\mathrm{r}|$. Le \textbf{pouvoir calorifique massique} est l'énergie libérée par kilogramme de combustible :
$\mathrm{PC}=\dfrac{|E_\mathrm{r}|}{M}$, en \unit{MJ.kg^{-1}}.
\end{definition}
\begin{propriete}[Quelques valeurs (eau formée à l'état gazeux)]
\par\smallskip\noindent\hfill\begin{tabular}{lccccc}
\toprule
Combustible & méthane & propane & butane & éthanol & dihydrogène \\
\midrule
$E_\mathrm{r}$ (\unit{kJ.mol^{-1}}) & $-802$ & $-2043$ & $-2657$ & $-1235$ & $-242$ \\
PC (\unit{MJ.kg^{-1}}) & 50,1 & 46,4 & 45,8 & 26,8 & 121 \\
\bottomrule
\end{tabular}\hfill\null\par\smallskip
\end{propriete}
\essSection{4}{Énergies de liaison}
\begin{definition}[Énergie de liaison]
L'\textbf{énergie de liaison} $D_{\mathrm{A-B}}$ est l'énergie qu'il faut fournir pour rompre une mole de liaisons A--B, en phase
gazeuse. En phase gazeuse, une réaction rompt des liaisons, ce qui consomme de l'énergie, et en forme de nouvelles, ce
qui en libère :
\[E_\mathrm{r}\approx\sum D(\text{liaisons rompues})-\sum D(\text{liaisons formées}).\]
\end{definition}
\begin{propriete}[Énergies de liaison moyennes (\unit{kJ.mol^{-1}})]
\par\smallskip\noindent\hfill\begin{tabular}{cccccccc}
\toprule
C--H & C--C & C--O & C=O (dans $\ce{CO2}$) & O=O & O--H & H--H \\
\midrule
415 & 345 & 358 & 804 & 498 & 463 & 436 \\
\bottomrule
\end{tabular}\hfill\null\par\smallskip
\end{propriete}
\schema{10-energie}
\begin{python}
D = {'C-H': 415, 'C-C': 345, 'C=O': 804, 'O=O': 498, 'O-H': 463}   # kJ/mol
rompues = {'C-H': 4, 'O=O': 2}           # CH4 + 2 O2
formees = {'C=O': 2, 'O-H': 4}           # CO2 + 2 H2O

E_rompues = sum(D[l] * n for l, n in rompues.items())
E_formees = sum(D[l] * n for l, n in formees.items())
Er = E_rompues - E_formees
print('Er =', Er, 'kJ/mol')
\end{python}
\begin{console}
Er = -804 kJ/mol
\end{console}
\noindent\capacite{Capacité numérique} Le programme associe à chaque liaison son énergie, puis fait la somme des énergies des liaisons
rompues et celle des liaisons formées. Le résultat, \qty{-804}{kJ.mol^{-1}}, est très proche de la valeur des tables, \qty{-802}{kJ.mol^{-1}}.
\essSection{5}{Combustions et enjeux de société}
\begin{propriete}[Usages, risques, perspectives]
Les combustions chauffent les logements, propulsent les véhicules et produisent une partie de l'électricité. Elles
comportent des risques : incendies, explosions, intoxications au monoxyde de carbone. Le dioxyde de carbone qu'elles
rejettent renforce l'effet de serre. Les axes d'étude actuels visent à réduire ces émissions : biocarburants, biogaz issu
de la méthanisation des déchets, dihydrogène produit par électrolyse, captage du dioxyde de carbone, et sobriété.
\end{propriete}
\schema{10-co2}
\begin{attention}
Une énergie molaire de réaction se rapporte à l'équation écrite : si l'on double tous les nombres stœchiométriques,
elle double. Les énergies de liaison sont des valeurs moyennes, valables en phase gazeuse : le calcul donne une estimation.
Ne pas confondre l'énergie libérée par mole et par kilogramme : le dihydrogène libère peu d'énergie par mole, mais
énormément par kilogramme.
\end{attention}
\begin{experience}[Estimer un pouvoir calorifique]
On chauffe une canette contenant \qty{200}{g} d'eau avec une lampe à éthanol, en mesurant l'élévation de température de l'eau
et la masse d'éthanol brûlé. L'énergie reçue par l'eau vaut $Q=m_\text{eau}\,c\,\Delta\theta$, avec $c=\qty{4,18}{kJ.kg^{-1}.\degreeCelsius^{-1}}$. Le pouvoir
calorifique estimé est très inférieur à la valeur des tables : une grande partie de l'énergie chauffe l'air ambiant.
\end{experience}
\begin{rediger}
Pour un bilan d'énergie : écrire l'équation ajustée, calculer la quantité de combustible, puis $Q=n\,|E_\mathrm{r}|$ ou $Q=m\times\mathrm{PC}$, en
précisant les unités. Avec les énergies de liaison, dresser la liste complète des liaisons rompues et formées, à l'aide des
schémas de Lewis.
\end{rediger}
\exercices

\rubrique{Découverte}
\begin{blocLivre}{ex-10-1}
\begin{exercice}[1]{Écrire des équations de combustion}
Écrire l'équation de la combustion complète de chaque espèce.
\begin{questions}
\item Le méthane $\ce{CH4}$ et le propane $\ce{C3H8}$.
\item Le butane $\ce{C4H10}$.
\item L'éthanol $\ce{C2H6O}$ et le méthanol $\ce{CH4O}$.
\item Vérifier, sur le propane et l'éthanol, les formules générales de l'essentiel.
\end{questions}
\end{exercice}
\end{blocLivre}
\begin{blocLivre}{ex-10-2}
\begin{exercice}[1]{Une combustion est une oxydoréduction}
On étudie la combustion du méthane, mettant en jeu les couples $\ce{CO2}/\ce{CH4}$ et $\ce{O2}/\ce{H2O}$.
\begin{questions}
\item Vérifier que la demi-équation $\ce{CO2 + 8H+ + 8e- = CH4 + 2H2O}$ est ajustée.
\item Écrire la demi-équation du couple $\ce{O2}/\ce{H2O}$ en milieu acide.
\item En déduire l'équation de la combustion du méthane.
\item Quel est le réducteur ? l'oxydant ? Que cède le méthane au dioxygène ?
\end{questions}
\end{exercice}
\end{blocLivre}
\begin{blocLivre}{ex-10-3}
\begin{exercice}[1]{Pouvoir calorifique}
\begin{donnees}
$E_\mathrm{r}(\text{méthane})=\qty{-802}{kJ.mol^{-1}}$ ; $E_\mathrm{r}(\text{propane})=\qty{-2043}{kJ.mol^{-1}}$ ; H 1,0 ; C 12,0 (\unit{g.mol^{-1}}).
\end{donnees}
\begin{questions}
\item Calculer les masses molaires du méthane et du propane.
\item Calculer leurs pouvoirs calorifiques massiques en \unit{MJ.kg^{-1}}.
\item Quel combustible libère le plus d'énergie par mole ? par kilogramme ?
\item Quelle énergie libère la combustion de \qty{1,0}{kg} de méthane ?
\end{questions}
\end{exercice}
\end{blocLivre}
\begin{blocLivre}{ex-10-4}
\begin{exercice}[1]{Chauffer de l'eau au butane}
On veut porter \qty{1,0}{L} d'eau de \qty{20}{\degreeCelsius} à \qty{100}{\degreeCelsius} avec un réchaud à butane. L'énergie nécessaire vaut $Q=m\,c\,\Delta\theta$,
avec $c=\qty{4,18}{kJ.kg^{-1}.\degreeCelsius^{-1}}$. Le pouvoir calorifique du butane vaut \qty{45,8}{MJ.kg^{-1}}.
\begin{questions}
\item Calculer l'énergie nécessaire.
\item Quelle masse de butane faudrait-il brûler si toute l'énergie libérée chauffait l'eau ?
\item En réalité, la moitié de l'énergie seulement chauffe l'eau. Quelle masse de butane faut-il ?
\item Où passe l'énergie « perdue » ?
\end{questions}
\end{exercice}
\end{blocLivre}
\begin{blocLivre}{ex-10-5}
\begin{exercice}[1]{Énergies de liaison et dihydrogène}
On étudie la réaction $\ce{2H2(g) + O2(g) -> 2H2O(g)}$, avec les énergies de liaison de l'essentiel.
\begin{questions}
\item Représenter les schémas de Lewis des réactifs et du produit.
\item Quelles liaisons sont rompues ? formées ? Combien de chaque ?
\item Calculer l'énergie de réaction pour l'équation écrite.
\item En déduire l'énergie libérée par mole de dihydrogène et la comparer à la valeur des tables, \qty{242}{kJ.mol^{-1}}.
\end{questions}
\end{exercice}
\end{blocLivre}
\rubrique{Entraînement}
\begin{blocLivre}{ex-10-6}
\begin{exercice}[2]{Énergies de liaison en Python}
On reprend le programme de l'essentiel.
\begin{questions}
\item \capacite{Capacité numérique}Pour la combustion du propane, $\ce{C3H8 + 5O2 -> 3CO2 + 4H2O}$, dénombrer les liaisons rompues et formées et écrire
les dictionnaires \code{rompues} et \code{formees}.
\item Prévoir l'affichage et comparer à la valeur des tables, \qty{-2043}{kJ.mol^{-1}}.
\item On écrit les dictionnaires pour l'équation $\ce{2C4H10 + 13O2 -> 8CO2 + 10H2O}$. Le programme affiche \code{Er = -5280 kJ/mol}. Quelle énergie
libère une mole de butane ? Comparer à \qty{-2657}{kJ.mol^{-1}}.
\item Pourquoi les résultats ne coïncident-ils pas exactement avec les valeurs des tables ?
\end{questions}
\end{exercice}
\end{blocLivre}
\begin{blocLivre}{ex-10-7}
\begin{exercice}[2]{Mesurer un pouvoir calorifique}
\schema{10-ex07}
On chauffe \qty{200}{g} d'eau ; sa température passe de \qty{18,0}{\degreeCelsius} à \qty{38,0}{\degreeCelsius}, tandis que la lampe perd \qty{1,20}{g} d'éthanol.
\begin{donnees}
$c=\qty{4,18}{kJ.kg^{-1}.\degreeCelsius^{-1}}$ ; pouvoir calorifique de l'éthanol : \qty{26,8}{MJ.kg^{-1}}.
\end{donnees}
\begin{questions}
\item Calculer l'énergie reçue par l'eau.
\item En déduire le pouvoir calorifique expérimental de l'éthanol, en \unit{kJ.g^{-1}}.
\item Comparer à la valeur des tables. Quelle part de l'énergie libérée a chauffé l'eau ?
\item Proposer deux améliorations du dispositif.
\end{questions}
\end{exercice}
\end{blocLivre}
\begin{blocLivre}{ex-10-8}
\begin{exercice}[2]{Se chauffer au gaz}
Une maison consomme \qty{15000}{kWh} de gaz naturel par an, assimilé à du méthane. On rappelle que $\qty{1}{kWh}=\qty{3,6}{MJ}$.
\begin{donnees}
Pouvoir calorifique du méthane : \qty{50,1}{MJ.kg^{-1}} ; $M(\ce{CH4})=\qty{16,0}{g.mol^{-1}}$ ; $M(\ce{CO2})=\qty{44,0}{g.mol^{-1}}$ ; $V_\mathrm{m}=\qty{24,0}{L.mol^{-1}}$.
\end{donnees}
\begin{questions}
\item Convertir la consommation annuelle en mégajoules.
\item Calculer la masse, puis la quantité de méthane brûlé en un an.
\item Calculer le volume de gaz consommé à \qty{20}{\degreeCelsius}.
\item Calculer la masse de dioxyde de carbone rejeté en un an.
\end{questions}
\end{exercice}
\end{blocLivre}
\begin{blocLivre}{ex-10-9}
\begin{exercice}[2]{Une combustion incomplète}
Une chaudière mal réglée, dans une pièce mal aérée, ne reçoit pas assez de dioxygène.
\begin{questions}
\item Écrire l'équation de la combustion du méthane produisant du monoxyde de carbone et de l'eau.
\item Écrire l'équation de la combustion produisant du carbone (suie) et de l'eau.
\item Pourquoi le monoxyde de carbone est-il particulièrement dangereux (chapitre~6) ?
\item Citer deux mesures de prévention.
\end{questions}
\end{exercice}
\end{blocLivre}
\begin{blocLivre}{ex-10-10}
\begin{exercice}[2]{Bois ou fioul ?}
Un foyer a besoin de \qty{10000}{kWh} de chaleur par an. Le bois sec a un pouvoir calorifique d'environ \qty{15}{MJ.kg^{-1}}, le fioul de
\qty{42,6}{MJ.kg^{-1}} (masse volumique \qty{0,84}{kg.L^{-1}}).
\begin{questions}
\item Calculer la masse de bois nécessaire.
\item Calculer la masse, puis le volume de fioul nécessaire.
\item Le bois rejette du dioxyde de carbone en brûlant. Pourquoi le considère-t-on pourtant comme une énergie
renouvelable ?
\item Quel inconvénient de la combustion du bois faut-il surveiller pour la qualité de l'air ?
\end{questions}
\end{exercice}
\end{blocLivre}
\begin{blocLivre}{ex-10-11}
\begin{exercice}[2]{Le dihydrogène, carburant d'avenir ?}
\begin{donnees}
Pouvoir calorifique : dihydrogène \qty{121}{MJ.kg^{-1}}, octane \qty{44,5}{MJ.kg^{-1}}. $V_\mathrm{m}=\qty{24,0}{L.mol^{-1}}$ à \qty{20}{\degreeCelsius} sous \qty{1}{bar} ; à
\qty{700}{bar}, la masse volumique du dihydrogène vaut environ \qty{40}{g.L^{-1}}.
\end{donnees}
\begin{questions}
\item Écrire l'équation de la combustion du dihydrogène. Quel est son principal avantage pour l'environnement ?
\item Quelle masse d'octane libère autant d'énergie que \qty{1}{kg} de dihydrogène ?
\item Quel volume occupe \qty{1}{kg} de dihydrogène à \qty{20}{\degreeCelsius} sous \qty{1}{bar} ? sous \qty{700}{bar} ?
\item Le dihydrogène est surtout produit aujourd'hui à partir de méthane. À quelle condition est-il un carburant « propre » ?
\end{questions}
\end{exercice}
\end{blocLivre}
\rubrique{Objectif DS}
\begin{blocLivre}{ex-10-12}
\begin{exercice}[3]{Réchaud à éthanol ou à butane ?}
\begin{docu}{1}{Deux réchauds de randonnée}
Les randonneurs utilisent des réchauds à cartouche de butane ou des réchauds à alcool (éthanol). L'éthanol peut être
produit à partir de betterave ou de canne à sucre.
\end{docu}
\begin{docu}{2}{Données}
$E_\mathrm{r}(\text{butane})=\qty{-2657}{kJ.mol^{-1}}$ ; $E_\mathrm{r}(\text{éthanol})=\qty{-1235}{kJ.mol^{-1}}$ ; pouvoirs calorifiques : butane \qty{45,8}{MJ.kg^{-1}}, éthanol
\qty{26,8}{MJ.kg^{-1}} ; $M(\ce{CO2})=\qty{44,0}{g.mol^{-1}}$.
\end{docu}
\begin{questions}
\item Écrire les équations de combustion complète du butane et de l'éthanol.
\item Quelle masse de chaque combustible faut-il emporter pour disposer de \qty{1}{MJ} ?
\item Calculer la masse de dioxyde de carbone rejetée par mégajoule pour chaque combustible.
\item \initiative Lequel choisir ? Argumenter, en tenant compte de l'origine du carbone de chaque combustible.
\end{questions}
\end{exercice}
\end{blocLivre}
\begin{blocLivre}{ex-10-13}
\begin{exercice}[3]{Le bilan énergétique d'une voiture}
\begin{docu}{1}{La consommation}
Une voiture consomme \qty{6,0}{L} d'essence aux \qty{100}{km}. On assimile l'essence à l'octane : masse volumique \qty{0,75}{kg.L^{-1}},
$M=\qty{114,0}{g.mol^{-1}}$, $E_\mathrm{r}=\qty{-5074}{kJ.mol^{-1}}$.
\end{docu}
\begin{docu}{2}{Le moteur}
Le moteur ne convertit en mouvement qu'environ \qty{35}{\percent} de l'énergie libérée par la combustion ; le reste est dissipé sous
forme de chaleur.
\end{docu}
\begin{questions}
\item Calculer la masse et la quantité d'octane brûlé par kilomètre.
\item Calculer l'énergie libérée par kilomètre, puis l'énergie utile au mouvement.
\item Calculer la masse de dioxyde de carbone émise par kilomètre.
\item \initiative Une voiture électrique consomme environ \qty{15}{kWh} aux \qty{100}{km}. Comparer à l'énergie chimique consommée par la
voiture à essence sur la même distance.
\end{questions}
\end{exercice}
\end{blocLivre}
\begin{blocLivre}{ex-10-14}
\begin{exercice}[3]{Le biogaz d'une ferme}
\begin{docu}{1}{La méthanisation}
Dans un méthaniseur, des bactéries décomposent les fumiers et les déchets agricoles en l'absence de dioxygène : il se
forme du biogaz, contenant environ \qty{60}{\percent} de méthane et \qty{40}{\percent} de dioxyde de carbone en volume.
\end{docu}
\begin{docu}{2}{La production}
Une exploitation agricole produit \qty{300}{m^3} de biogaz par jour. Un foyer consomme environ \qty{15}{MWh} d'énergie par an pour son
chauffage et son eau chaude.
\end{docu}
\begin{donnees}
$V_\mathrm{m}=\qty{24,0}{L.mol^{-1}}$ ; $E_\mathrm{r}(\text{méthane})=\qty{-802}{kJ.mol^{-1}}$ ; $\qty{1}{kWh}=\qty{3,6}{MJ}$.
\end{donnees}
\begin{questions}
\item Calculer la quantité de méthane contenue dans \qty{1}{m^3} de biogaz à \qty{20}{\degreeCelsius}.
\item Calculer l'énergie libérée par la combustion de \qty{1}{m^3} de biogaz, en mégajoules puis en kilowattheures.
\item Calculer l'énergie produite en un an par l'exploitation.
\item \initiative Combien de foyers ce biogaz pourrait-il chauffer ? Pourquoi parle-t-on d'une énergie renouvelable ?
\end{questions}
\end{exercice}
\end{blocLivre}
\begin{blocLivre}{ex-10-15}
\begin{exercice}[3]{La puissance d'une bougie}
\begin{docu}{1}{Une bougie}
Une bougie est formée de paraffine, que l'on assimile à l'alcane $\ce{C25H52}$, de pouvoir calorifique \qty{42}{MJ.kg^{-1}}. En une heure,
la bougie perd \qty{5,0}{g}.
\end{docu}
\begin{docu}{2}{Une comparaison}
Une ampoule à DEL éclaire autant qu'une bougie avec une puissance d'environ \qty{0,1}{W}.
\end{docu}
\begin{questions}
\item Écrire l'équation de la combustion complète de la paraffine.
\item Calculer l'énergie libérée en une heure.
\item \initiative Calculer la puissance thermique de la bougie, en watts, et la comparer à celle de l'ampoule.
\item Calculer la masse de dioxyde de carbone rejetée en une heure.
\end{questions}
\end{exercice}
\end{blocLivre}
